\subsection{THETA3-A15. The literal finite transform, its derivative matrix, and rank-one sampling}
\label{A15}

We first compute the actual finite transform for the whole allowed
array.  We then prove a positive first-derivative theorem for its
single-incidence Fourier model.  The final paragraph specifies the
additional tensor-product test required when several $T$-primes occur.

\paragraph{The exact operator and first derivative.}
Use the actual data of
\texttt{Development/AbelianVarietyTheory/Proofs/DecoratedLatticeProofs/}
\texttt{DatumReadingSplitPlacesEqClassicalTheta.lean},
\texttt{reads\_datum\_eq\_thC}.  In particular the source has a valid
\texttt{StdPair}, the stipulated half idempotent, self-dual lattices,
vacuum and \texttt{IsThetaAt} identity, and the required support and
periodicity of the actual array $\Xi$.  Set
\[
 H=\mathrm{HV}(\mathrm{sc}),\quad
 \alpha=\mathrm{alV}(d,\mathrm{sc}),\quad
 F=\mathrm{xiN}(f_0).
\]
Here $H$ is linear in its first argument and conjugate-linear in its
second.  The exact definition, including its sign, is
\texttt{Formalism/LocalToGlobal/TorsionIntegrality/QuadraticSemicharacters/}
\texttt{IsSemiChar.lean}, lines 9--24:
\[
 U_xF(z)=F(z+x)
       \exp\bigl(-\pi H(z,x)-\tfrac\pi2H(x,x)\bigr).
 \tag{A15a}
\]
Write $S_X,S_A,S_B$ for the three finite representative sets in the
source proof, and put
\[
 b(y,s)=\exp(-2\pi i\,\operatorname{Im}H(s,y)),
 \qquad
 K_\Xi(x,y)=\frac1{|S_A|}\sum_{s\in S_A}
          \Phi(s-x)b(y,s)^{-1},
 \tag{A15b}
\]
where $\Phi=\mathrm{onV}(\Xi)$.  Then the literal conclusion of that
proof is
\[
 \Theta_\Xi(z)=
  \sum_{x\in S_X,\ y\in S_B}
      \alpha(x)\alpha(y)K_\Xi(x,y)U_xU_yF(z).
 \tag{A15c}
\]
The definitions are
\texttt{Formalism/LocalToGlobal/KernelLattices/XiIndCClass.lean},
\texttt{brk}, \texttt{hatS}, \texttt{xiL}, lines 18--28,
and \texttt{thC}, line 239.  The actual \texttt{thC} uses volume
$1$.  If instead one uses the separately defined \texttt{xiC},
its displayed positive volume contributes the additional factor
$\sqrt{\mathrm{vol}}$; it is not absorbed in (A15c).

For a tangent vector $u$, direct differentiation of (A15a) gives
\[
 \begin{split}
 d(U_xU_yF)_0(u)
  ={}&\exp\bigl(-\pi H(x,y)
              -\tfrac\pi2H(x,x)-\tfrac\pi2H(y,y)\bigr)\,V_u(x+y),\\
 V_u(w)={}&dF_w(u)-\pi H(u,w)F(w).
 \end{split}
 \tag{A15d}
\]
Indeed $H(z,x)+H(z,y)=H(z,x+y)$, and the constant term in
the two exponential factors is precisely the exponent displayed.
This proves (A15d) without a convention-dependent appeal to a
Heisenberg commutation rule.  Equivalently,
$U_xU_y=\exp(-\pi i\operatorname{Im}H(x,y))U_{x+y}$.

Fix tangent vectors $u_1,\ldots,u_g$.  The complete first-derivative
matrix for any finite list of actual arrays $\Xi_j$ is therefore
\[
 J_{rj}=\frac1{|S_A|}
  \sum_{x,y,s}\alpha(x)\alpha(y)b(y,s)^{-1}
       \Phi_j(s-x)
       e^{-\pi H(x,y)-\pi H(x,x)/2-\pi H(y,y)/2}
       V_{u_r}(x+y).
 \tag{A15e}
\]
Every index and scalar in this formula comes from the stated finite
transform.  In particular $\Xi_j(0)=0$ does not set $F(0)$ or the
translated vacuum values to zero.  A prescribed full-rank support
condition or right-$\mathrm{SL}$ covariance is retained inside
$\Phi_j$ throughout this formula.

There is a useful rigorous error certificate for (A15e).  Suppose
computed values of $F(w)$ and $dF_w(u_r)$ have respective error bounds
$\epsilon_0(w)$ and $\epsilon_r(w)$.  Since the semicharacters and
commutator factors have modulus one, the resulting error in $J_{rj}$
is at most
\[
 \frac1{|S_A|}\sum_{x,y,s}|\Phi_j(s-x)|
 e^{-\pi H(x+y,x+y)/2}
 \bigl(\epsilon_r(x+y)+\pi|H(u_r,x+y)|\epsilon_0(x+y)\bigr).
 \tag{A15f}
\]
The real part of the exponent in (A15e) is
$-\pi H(x+y,x+y)/2$, which proves this estimate.  The theta and
derivative tail bounds of Q2 and A12 can supply the errors here.
Thus a numerical nonzero-minor certificate need not omit any of the
finite half-pairing phases.

\paragraph{The exact restricted Fourier rank test.}
Fix all non-$T$ data, including an actually permitted primary tuple
from A14 and its allowed local indices.  At a $T$-prime $v$, let
$M_v=\operatorname{Mat}_{b\times a}(\kappa_v)$ in the source's residue
frames.  Integrality is included in the remaining factor, so on the
finite window in (A15e) the actual array has the form
\[
 C_j(t)\prod_{v\in\mathcal T}I_{v_v,u_v}(\operatorname{red}_v t).
 \tag{A15g}
\]
Here $C_j$ is the \emph{actual} product of other factors, including
the indicated finite auxiliary sums.  It is not replaced by a
constant.  Collecting all terms with a fixed tuple of reductions
defines a completely explicit weighted kernel
$\mathcal K_{rj}$ on $\prod_v M_v$, so that (A15e) becomes
\[
 J_{r,(j,I)}=\sum_{X\in\prod_v M_v}
                    \mathcal K_{rj}(X)\prod_v I_v(X_v).
 \tag{A15h}
\]
The collection operation simply sums the coefficient of
$C_j(t)$ in (A15e) over all triples $(x,y,s)$ with $t=s-x$ and
the specified reductions.  Terms off the required integral domain
are zero.  Its definition therefore involves only the literal
finite window, not an infinite summation or a changed support.

Choose nontrivial additive residue characters.  For a matrix $A$
put $e_A(X)=\psi(\operatorname{Tr}(A^tX))$.  A13 implies that the
single-prime incidence space is also exactly
\[
 \operatorname{span}\{e_A-1:\operatorname{rank}A=1\}.
 \tag{A15i}
\]
For completeness, if $A=v^tt$, then
\[
 e_A(X)-1
       =\sum_{u\ne0}\bigl(\psi(\textstyle\sum_j u_jt_j)-1\bigr)I_{v,u}(X).
 \tag{A15j}
\]
The value at $vX=0$ is zero, so the missing $u=0$ term causes no
error.  Conversely the Fourier support of each incidence function
lies in $\{v^tt:t\in\kappa^a\}$ and its value at zero is zero,
as proved in A13.  Fourier independence then proves (A15i).
The coefficient sum in (A15j) is at most $2(q^a-1)$.

Use the positive-sign finite transform
\[
 \widehat{\mathcal K}_{rj}(A_1,\ldots,A_f)
   =\sum_X\mathcal K_{rj}(X)\prod_v e_{A_v}(X_v).
\]
For a derivative covector $\ell$, put
$\mathcal K_{\ell j}=\sum_r\ell_r\mathcal K_{rj}$.
The actual incidence derivative matrix has full row rank if and
only if there is no nonzero $\ell$ for which
\[
 \sum_{S\subset\{1,\ldots,f\}}(-1)^{f-|S|}
       \widehat{\mathcal K}_{\ell j}(A^S)=0
 \quad\text{for all }j\text{ and all rank-one }A_v.
 \tag{A15k}
\]
Here $A^S_v=A_v$ for $v\in S$ and $A^S_v=0$ otherwise.
Indeed the left side is exactly the pairing in (A15h) with
$\prod_v(e_{A_v}-1)$, and these products span the tensor incidence
space.  This proves both directions of the criterion.  In the
single-prime case it says that
$\widehat{\mathcal K}_{\ell j}$ is constant on all matrices of rank
at most one.  For several primes it is the full mixed difference
(A15k), rather than the separate one-variable condition.
Equivalently, on this rank-one grid the transform is the alternating
sum of its proper coordinate faces.  The number of local basis
columns before multiplying the other factors is
\[
       \prod_{v\in\mathcal T}
       \frac{(q_v^a-1)(q_v^b-1)}{q_v-1}.
 \tag{A15l}
\]

When varying the primary twisting characters, the index $j$ runs
only over the quotient-character family actually constructed in
A14.  In particular the finite auxiliary sum in a \texttt{XiEntry}
is not silently replaced by one determinant character.  When passing
from raw decorated lattices to the archive's chosen class
representatives, include its stated equivariance scalar and its row
change of coordinates in (A15e).  The criterion then tests the
resulting actual matrix.  No enlarged list of class representatives
is assumed.

\paragraph{Normalization and a finite algebraic certificate.}
The literal differential convention is especially simple at degree
one.  In
\texttt{Formalism/LocalToGlobal/KernelLattices/ThDistC.lean},
lines 54--62, \texttt{jetC} is the undivided derivative multiplied by
$c_D(\tau)$ in each differentiated coordinate.  For
$r=(\tau,i,j)$ use
$u_r=c_D(\tau)\,\mathrm{coordX}(r)$ in (A15e).
Then the result is exactly its degree-one \texttt{jetC}; there is
no factorial.  A different rational Schr\"odinger chart requires
the actual row Jacobian retained in A12.

For fixed source and target lattices, divide the entries by the
common factor $\chi_F(N,N')$ specified in A13.  The fixed-$J$
arithmetic-jet construction there, with degree one included, puts
all entries for the finite actual family in one number field after
a further \texttt{Enl}.  This invokes the literal normalized-jet
theorem at its stated \texttt{AmbIn.Ok} hypotheses; it does not
deduce algebraicity of unnormalized analytic derivatives.
Changing one theta section into a linear combination still uses
its section comparison factor once.  A determinant of independent
first-derivative columns uses the corresponding product of factors,
as in A12.

Thus a nonzero minor of this \emph{normalized} matrix gives an
algebraic inverse on its image.  If a computed square submatrix
$J_0$ has least singular value at least $s>0$, then a target
derivative vector $w$ can be obtained with coefficient sum at most
$\sqrt g\,\|w\|/s$.  This follows from
$\|J_0^{-1}w\|_1\leq\sqrt g\|J_0^{-1}w\|_2$.
For an approximate matrix replace $s$ by
$\sigma_{\min}(J_{0,\mathrm{approx}})-\|\mathrm{error}\|$, provided
this is positive; (A15f) bounds the error.
These are concrete finite certificates.  An algebraic inverse of
this derivative matrix is still not, without a cohomological
identification, an algebraic correspondence on the ball variety.

\paragraph{Faithfulness and modulation in the precise finite block.}
There is a useful positive nonvanishing criterion before any derivative
is taken.  Suppose that $\Lambda,A,B$ are commensurable full lattices,
that all three are self-dual for $\operatorname{Im}H$, and that
$\alpha$ is a semicharacter on each.  Retain the actual
\texttt{GoodLat} conditions: in particular
\[
 \Lambda=(\Lambda\cap A)+(\Lambda\cap B),\qquad
 a\in A,\ b\in B,\ a+b\in\Lambda\ \Longrightarrow\ a\in\Lambda.
 \tag{A15u}
\]
These are the first two lattice conditions in
\texttt{Formalism/LocalToGlobal/PairingIndexBounds/ModelLatticeMoves/}
\texttt{GoodLat.lean}, lines 25--29; its last two conditions assert
the corresponding nondegeneracy of the cross pairing.
Let $F_0\ne0$ satisfy $\alpha(\lambda)U_\lambda F_0=F_0$.
Then (A15c) is injective on the finite space of functions
\[
                 \Phi:A/(A\cap B)\longrightarrow\C.
 \tag{A15v}
\]
In this statement $S_X$ represents $A/(\Lambda\cap A)$,
$S_A$ represents $(\Lambda\cap A)/(A\cap B)$, and $S_B$
represents $B/(\Lambda\cap B)$, precisely as in the source.
The inclusion $A\cap B\subset\Lambda$ follows from (A15u): an
element of $A\cap B$ pairs integrally with both summands of
$\Lambda$, hence belongs to its self-dual lattice.

Here is a full proof of injectivity.  For every $\lambda\in\Lambda$,
the nonzero entire function $U_xU_yF_0$ is an eigenvector of
$\alpha(\lambda)U_\lambda$, with character
\[
             e^{-2\pi i\operatorname{Im}H(\lambda,x+y)}.
\]
These characters are distinct as $(x,y)$ runs over $S_X\times S_B$.
Equality would imply $x+y-x'-y'\in\Lambda$ by self-duality;
(A15u) then gives $x-x'\in\Lambda$, forcing $x=x'$ and $y=y'$
by the representative conditions.  Eigenvectors with distinct
characters of an abelian group are linearly independent.  For
example, given a finite relation and a chosen character, successive
operators of the form $\rho(\lambda)-\chi(\lambda)$ eliminate
each of the other characters while leaving that chosen term
nonzero.  Thus $\Theta_\Phi=0$ forces $K_\Phi(x,y)=0$ for every
pair.

The pairing $b(y,s)$ between
$(\Lambda\cap A)/(A\cap B)$ and $B/(\Lambda\cap B)$ is perfect.
Its right kernel is exactly $\Lambda\cap B$ by the last
\texttt{GoodLat} condition.  Its left kernel is $A\cap B$, since
an element of $\Lambda\cap A$ pairing integrally with all of $B$
belongs to $B$ by self-duality.  Both groups are finite, so
nondegeneracy on both sides makes their orders equal and gives a
perfect finite-character pairing.  Finite Fourier inversion in
(A15b) now gives $\Phi(s-x)=0$ for every $s,x$.
These elements represent all of $A/(A\cap B)$: choose $x$
representing $-t$ modulo $\Lambda\cap A$, then $s$ representing
$t+x$ modulo $A\cap B$.  This proves (A15v).

There is also an exact operator identity on this whole block,
including an already present array of other local factors:
\[
       \Theta_{\,b(y,\cdot)\Phi}
                =\alpha(y)U_y\Theta_\Phi,
                  \qquad y\in B.
 \tag{A15w}
\]
To verify it without dropping a phase, multiplication of $\Phi(s-x)$
by $b(y,s-x)$ changes $K_\Phi(x,z)$ to
$b(y,-x)K_\Phi(x,z-y)$.
The coefficient $K_\Phi(x,z)$ is periodic in
$z\in B$ modulo $\Lambda\cap B$, because $s\in\Lambda\cap A$.
On $B$ the operators $W_y=\alpha(y)U_y$ form a representation:
the semicharacter identity and
$U_xU_y=e^{-\pi i\operatorname{Im}H(x,y)}U_{x+y}$ give
$W_xW_y=W_{x+y}$, since $\operatorname{Im}H(B,B)\subset\Z$.
Also $W_\lambda F_0=F_0$ for $\lambda\in\Lambda\cap B$.
Reindexing $z-y$ in $S_B$ therefore moves $W_y$ to the left of
$W_zF_0$.  Commuting it past $U_x$ contributes $b(y,x)$,
which cancels $b(y,-x)$.  This gives exactly (A15w).
It is the statement called \texttt{XiModulate} in
\texttt{XiIndCClass.lean}, lines 98--102; the argument here proves
the finite identity directly, rather than treating that
\texttt{Prop} declaration as a proof.

Consequently, for any collection $y_1,\ldots,y_f\in B$,
\[
 \Theta_{\,\Phi\prod_i(b(y_i,\cdot)-1)}
      =(W_{y_f}-1)\cdots(W_{y_1}-1)\Theta_\Phi.
 \tag{A15x}
\]
This proves the ordered character-operator prerequisite in C37
whenever the literal input has been placed in this finite block.
The representatives need not be bounded; C37 retains the lattice
characters introduced by reducing them modulo the baseline period
lattice.  The positive-sign finite character $e_A$ in (A15i) may
label $-y$ rather than $y$ under the source's negative-sign
$b(y,t)=\exp(-2\pi i\operatorname{Im}H(t,y))$.
Use that actual perfect-pairing label.  Negation preserves the
rank-one sample and does not change any spanning assertion.

In particular a nonzero remaining array gives a nonzero baseline
\emph{when its membership in (A15v) and these self-duality and
vacuum hypotheses have been checked}.  This is more specific than
an assertion about an arbitrary theta summation map.  To use one
fixed baseline while a new split prime varies, the added local
factor must be the unramified integral vacuum at that prime.
In the finite Fourier formula its integral support sets the new
$x$-index to zero, and its constant value on the new $S_A$ makes
the new Fourier index zero.  Thus this factor acts as the identity.
Equivalently, the orthogonal prime components of the finite
Heisenberg group factor by the Chinese remainder theorem, and
the new factor is its vacuum vector.  This statement keeps the
remaining lattices, section normalization and local factors fixed;
it does not assert that an arbitrary changing archive datum has a
fixed baseline.  For fixed lattices, the new sufficiently large
primes avoid their finite nonstandard set, so the added local
factor of the volume is one.  Any pre-existing positive volume
scalar and every later class-representative or row-frame factor
remain explicit, as in (A15e) and the normalization paragraph.

\paragraph{The fixed covering lattice and the new primary samples.}
The baseline in the preceding paragraph need not be a section on
the original principally polarized torus.  Let $M$ be its original
$\OO_E$-period module, with complex image $\Lambda$, and choose a
positive integer $N_0$ such that
\[
 \Lambda_0=N_0\Lambda\subset A_{\rm old}\cap B_{\rm old}.
 \tag{A15y}
\]
Such an integer exists because the old lattices are commensurable.
For $\lambda\in\Lambda_0$, the character $b(\lambda,t)$ is one
on $A_{\rm old}$, by self-duality of $A_{\rm old}$.
Equation (A15w) therefore proves
$\alpha(\lambda)U_\lambda F_{\rm base}=F_{\rm base}$.
Thus the actual fixed positive theta line for C37 is its line on
$V/\Lambda_0$.  No descent of $F_{\rm base}$ to $V/\Lambda$
is being assumed.

Here is the precise new-prime compatibility for the half-lattices
used above.  Exclude primes dividing $N_0$, the old levels, the
different, and the finitely many nonstandard lattice indices.
At a new split rational prime $q$, the source definitions
\texttt{halfAO}, \texttt{halfBO} in
\texttt{XiIndCClass.lean}, lines 181--191, give the local pairs
$q^{-1}M,qM$ at each selected prime $\mathfrak p$ and the
opposite pair at $\overline{\mathfrak p}$.
The baseline has the integral vacuum at these places.
The $\overline{\mathfrak p}$-primary torsion group
\[
 \overline{\mathfrak p}^{-1}(N_0M)/(N_0M)
\]
on $V/\Lambda_0$ maps isomorphically to the corresponding
primary group on $V/\Lambda$.
Indeed, the kernel of $V/\Lambda_0\to V/\Lambda$ has order
prime to $q$, so multiplication by $q$ is invertible on it.
This proves injectivity on $q$-torsion, and a lift of a
$q$-torsion point can be adjusted by this kernel to become
$q$-torsion.  The map is $\OO_E$-linear, so it preserves the
individual prime components.

One can use these very primary points in (A15w).
Represent a point by $y\in\overline{\mathfrak p}^{-1}(N_0M)$.
By the Chinese remainder theorem for the finite module
$N_0M/qN_0M$, add an element of $N_0M$ so that $y$ belongs
to $qM$ at $\mathfrak p$ and at every other prime above $q$
except $\overline{\mathfrak p}$.  At the old primes its
component remains in $N_0M\subset A_{\rm old}\cap B_{\rm old}$.
The resulting representative belongs to $B_{\rm new}$.
On the baseline's integral new-prime support its character
$b(y,t)$ is precisely the residue trace character at
$\mathfrak p$, with the negative sign and scalar different
factor already recorded in (A15p); all other prime components
are trivial.  Perfect local trace duality proves this assertion
coordinate by coordinate.  With the stipulated source and target
frames the labels are matrices, and replacing $M$ by $N_0M$
only multiplies their local labels by a nonzero scalar.  Thus the
rank-one labels remain rank one.  The same construction applies
independently to every selected prime above $q$.

Consequently the torsion samples needed by C37 are samples on
the fixed torus $V/\Lambda_0$, with its fixed baseline section.
Their trace-dual Fourier test is (A15p), now for the dual of
$\Lambda_0$; a fixed nonzero dual element reduces nontrivially
away from finitely many primes.  The representatives in C37 may
be unbounded, but their changes of lift are elements of
$\Lambda_0$, as required there.  A change merely by an element
of the larger lattice $\Lambda$ is not declared harmless.
This argument proves compatibility under the stated literal
half-lattice, integral-vacuum and tensor-frame hypotheses; it
does not replace those hypotheses for an unrelated finite
transform.

\paragraph{A positive theorem for one incidence factor.}
Here is a nonvanishing result for the restricted rank-one frequencies
themselves.  Let $V/\Lambda$ be a complex torus with a positive theta
line, Hermitian form $H$, and semicharacter $\alpha$.  Let $F\ne0$
be one of its holomorphic theta sections in the entire-function
frame used in (A15a).  Set
\[
       v(y)=d_z[U_yF](0)\in V^*,\qquad a_0=v(0).
 \tag{A15m}
\]
Although $v$ depends on a unitary frame, its norm and the norm of
each constant covector pairing are well defined on the torus.
More precisely, $\alpha(\lambda)U_\lambda F=F$ implies
\[
 U_{y+\lambda}F
    =\alpha(\lambda)^{-1}
       e^{\pi i\operatorname{Im}H(y,\lambda)}U_yF.
 \tag{A15n}
\]
The same scalar multiplies every component of $v$, so those
components are continuous sections of the same nontrivial theta
line in a unitary frame.

This parameter-space statement does not make every $U_y$ an
automorphism of the original theta line on $V/\Lambda$.
For a point with $Ny\in\Lambda$, $U_yF$ and $F$ have the same
automorphy law on $N\Lambda$, so they are sections together on
$V/N\Lambda$ of the pullback of the original theta line under
the identity-linear covering
$\pi_N:V/N\Lambda\longrightarrow V/\Lambda$.
Indeed the extra multiplier on $N\lambda$ is
$e^{2\pi i\operatorname{Im}H(y,N\lambda)}=1$.
In this covering's $z$ coordinate the derivative of $\pi_N$ is
the identity, and (A15m) is its literal derivative.  If the cover
is instead identified with the original torus by $z=Nw$, its
line is $[N]^*L$ and its first derivative is multiplied by $N$.
Only the explicitly normalized factor $N^{-1}$ cancels that
change of coordinate.  The fixed positive line in (A15o) is
the parameter line of $v(y)$; the output columns for growing
torsion levels live on these stated pullback lines.

For any fixed Hermitian norm on $V^*$, put
\[
 \delta_F=\min_{\|\ell\|=1}
     \int_{V/\Lambda}
       \bigl(|\ell v(y)|-|\ell a_0|\bigr)^2\,dy,
 \tag{A15o}
\]
where Haar measure is a probability measure.  Then
$\delta_F>0$.  To prove this, if the integrand has zero integral
it vanishes everywhere by continuity.  If its constant norm is
positive, the corresponding nowhere-zero section trivializes the
theta line topologically, contradicting its nonzero positive
first Chern class.  This contradiction can be read directly from
curvature: a global unit frame would make its curvature exact,
whereas its positive curvature has a positive integral against
the $(g-1)$st power of a K\"ahler form.  If the constant is zero,
the directional form of (A15d) gives, for the nonzero direction
$u$ corresponding to $\ell$,
\[
             dF_y(u)=\pi H(u,y)F(y)\quad\hbox{for all }y.
\]
Apply the antiholomorphic derivative in direction $u$.  Holomorphy
of $F$ gives $0=\pi H(u,u)F(y)$, which contradicts $F\ne0$ and
positive definiteness.  Thus the integral in (A15o) is positive
for every unit $\ell$.  It is continuous in $\ell$, and the unit
sphere is compact, proving the positive minimum.

Suppose the finite torsion groups at a sequence of allowed CM split
primes are labelled, in tensor-compatible residue frames, by
$\operatorname{Mat}_{b\times a}(\kappa)$, with $q=|\kappa|\to\infty$.
Write $y_A$ for the point labelled by $A$.  The compatibility required
here is the actual trace-pairing compatibility: for every torus
character $\gamma$ there is a residue matrix $X_\gamma$ with
\[
        \gamma(y_A)=\psi(\operatorname{Tr}(A^tX_\gamma)),
 \tag{A15p}
\]
and every fixed nonzero $\gamma$ has $X_\gamma\ne0$ at all but
finitely many of the primes.  This is the usual CM lattice
construction: the rational trace dual is an $E$-lattice, and a
nonzero element of it has nonzero reduction away from finitely many
primes.  At good split primes the trace pairing is perfect.
Left and right changes of the stipulated source and target frames,
and the scalar different factor, preserve rank one.  An arbitrary
unrecorded linear relabelling of matrix coordinates is not allowed
in (A15p).

The probability measure on the $y_A$ with $\operatorname{rank}A=1$
converges to Haar measure.  Indeed every such matrix is $uv^t$ with
$u,v\ne0$, in exactly $q-1$ ways.  If $r=\operatorname{rank}X$, direct
summation gives
\[
 \frac{\sum_{u\ne0,\ v\ne0}\psi(u^tXv)}
            {(q^b-1)(q^a-1)}
  =\frac{q^{a+b-r}-q^a-q^b+1}{(q^b-1)(q^a-1)}.
 \tag{A15q}
\]
For fixed $u$ the inner sum is $q^a-1$ if $X^tu=0$ and $-1$
otherwise; there are $q^{b-r}-1$ nonzero $u$ of the first kind.
This proves the formula.  For $r\ge1$ its absolute value is at
most $1/(q-1)$: the largest positive numerator occurs at $r=1$,
and the negative magnitude is at most
$\min(q^a,q^b)-1$, both bounded by the indicated denominator
divided by $q-1$.  Equations (A15p)--(A15q) give convergence on
each torus character, hence on trigonometric polynomials, and then
on continuous functions by uniform approximation.

For \emph{any} choices of phases $|\eta_A|=1$, the columns
\[
                 \eta_Av(y_A)-a_0,
                 \qquad \operatorname{rank}A=1,
 \tag{A15r}
\]
span $V^*$ for every sufficiently large such prime.  More precisely,
if $R_q=(q^a-1)(q^b-1)/(q-1)$ is their number, their matrix $J_q$
satisfies
\[
             J_qJ_q^*\ \ge\ R_q\delta_F/2\;\mathrm{Id}.
 \tag{A15s}
\]
The triangle inequality gives
$|\eta_A\ell v(y_A)-\ell a_0|
 \ge\bigl||\ell v(y_A)|-|\ell a_0|\bigr|$.
Equidistribution is uniform on the compact family of continuous
functions in (A15o): use a finite uniform approximation net in
$\ell$ and then the proved convergence for that finite net.
Their discrete integrals are therefore at least $\delta_F/2$,
uniformly in unit $\ell$ and in every phase.  This proves (A15s).
The minimum-norm right inverse gives coefficient sum at most
$\sqrt{2/\delta_F}\,\|w\|$ for a desired derivative vector $w$.
To express each difference by the literal incidence indicators,
multiply this bound by at most $2(q^a-1)$ using (A15j).

There is also a finite quantitative threshold.  Choose coordinates
on the real torus of dimension $d=2g$, and let $P_M$ be the product
Fej\'er kernel of frequency range $[-M,M]^d$.  Choose a uniform
bound $B\ge\sup_y\|v(y)\|$ and a bound $L$ for its covariant first
derivative.  The functions in (A15o) have bound $4B^2$ and
Lipschitz constant at most $4BL$.  Put
$\epsilon_M=\int\operatorname{dist}(t,0)P_M(t)\,dt$.
After excluding the finitely many primes where one of the nonzero
frequencies in this range reduces to zero, (A15q) shows that the
discrete average in (A15o) is at least
\[
 \delta_F-4BL\epsilon_M
       -\frac{4B^2((2M+1)^d-1)}{q-1}.
 \tag{A15t}
\]
Convolution by $P_M$ preserves the Haar average, approximates every
function in the family within $4BL\epsilon_M$, and has the
indicated finite Fourier support; each Fourier coefficient has
absolute value at most $4B^2$.  This proves (A15t), including its
constants.  Taking the last two terms at most $\delta_F/4$ each
is an explicit sufficient test for (A15s).  The constants concern
this fixed theta section; they do not bound an arbitrary varying
section or an unnormalized ball period.

\paragraph{Application to several factors.}
In the single-incidence Fourier model the exact finite transform
sends $e_A$ to a unit scalar times $U_{y_A}F$, and sends $1$ to
$F$.  Thus (A15r) is exactly its first-derivative column.  The
theorem proves that low Fourier rank alone is not an obstruction
to first-derivative spanning, while retaining every unit phase.
For the literal full \texttt{xiAt}, however, all the $T$-factors,
the other local factors, and the allowed primary tuple list remain
those in (A15g).  Several incidence factors give the alternating
sum (A15k).  C37 proves a simultaneous theorem for these ordered
differences, at one growing rational prime, even with unbounded
half-lattice representatives.  The exact identity (A15x) supplies
its operator comparison in the stated finite block.  The
injectivity proof supplies a nonzero baseline from a nonzero
array in that block, and the unramified-factor argument explains
when this is one fixed baseline as the new prime grows.
Together these results give full first-derivative rank whenever
those finite-block, fixed-data, and tensor-compatible residue
conditions are met for the sequence of actual inputs.

For a particular archive family one must still verify these
conditions for its selected non-$T$ local functions and permitted
prime choices; they are not consequences merely of a nonzero raw
array or a numerical signature match.  The exact finite
nonzero-minor test, normalized arithmetic field, and error bounds
remain (A15e)--(A15l) and (A15f).  In particular the literal
$v_T(u)$ has several places over the same rational $T(u)$;
C37 addresses that simultaneous situation without assuming
independent successive prime choices.  All these comparisons
retain the actual allowed array space of A13--A14.

In the requested Grothendieck framework these are comparisons of
finite theta values and their normalized derivatives, as in
\emph{cote44}, scans 7--8.  Turning their inverse into a map of
homological motives requires the separate cohomological
identification; it is not supplied by finite-dimensionality or by
the algebraic Lefschetz inverses discussed in \emph{cote16},
scans 43--46.
